\section*{Comentarios y conclusiones}

El presente documento consolida la descripción técnica de alto nivel y el conjunto de requerimientos de producto y sistema del collar inteligente para cerca virtual (DR-VF-01). Se establecieron las funciones esenciales que debe cumplir el dispositivo, así como las restricciones de diseño térmico, energético y de aislamiento galvánico.

Este documento establece la línea base sobre la cual se realizarán el diseño esquemático detallado en Altium Designer / EasyEDA, la fabricación de la placa de circuito impreso de prueba (\emph{Carrier Board}) y el desarrollo del firmware en Zephyr RTOS. Entre los aspectos más relevantes del diseño destacan:
\begin{itemize}[leftmargin=*]
    \item La selección del regulador solar \textbf{CN3791 de 6V MPPT}, el cual maximiza la recarga de la batería LiPo sin elevar la complejidad ni el costo de la tarjeta de circuito impreso.
    \item La utilización de \textbf{LoRa SX1262} como red primaria de comunicaciones, resolviendo la falta de cobertura celular en zonas rurales extensas sin generar costos operativos por suscripción.
    \item La reserva estricta de recursos de hardware en el ESP32-S3 (puerto UART2 y rieles de corriente $V_{\text{BAT}}$), garantizando que el diseño sea escalable hacia la integración del módem celular \textbf{Quectel BG95-M3} durante la fase de tesis.
\end{itemize}

En conclusión, esta especificación gestiona adecuadamente las expectativas del proyecto al diferenciar con precisión lo que se entrega en el semestre actual y las metas de desarrollo a largo plazo, enfocando el trabajo de ingeniería en un prototipo funcional, seguro y comercialmente viable.